\documentclass[aspectratio=169,10pt]{beamer}

% Shared Beamer setup for Fall 2026 course slides.
\mode<presentation>{
  \usetheme{Madrid}
  \usecolortheme{default}
}

\definecolor{CalPolyGreen}{HTML}{154734}
\definecolor{CalPolyGold}{HTML}{BD8B13}
\definecolor{DeckInk}{HTML}{1F2933}
\definecolor{DeckMuted}{HTML}{5F6B7A}

\setbeamercolor{structure}{fg=CalPolyGreen}
\setbeamercolor{palette primary}{bg=CalPolyGreen,fg=white}
\setbeamercolor{palette secondary}{bg=DeckInk,fg=white}
\setbeamercolor{palette tertiary}{bg=CalPolyGold,fg=DeckInk}
\setbeamercolor{frametitle}{bg=CalPolyGreen,fg=white}
\setbeamercolor{title}{fg=CalPolyGreen}
\setbeamercolor{normal text}{fg=DeckInk,bg=white}
\setbeamercolor{block title}{bg=CalPolyGreen,fg=white}
\setbeamercolor{block body}{bg=black!3,fg=DeckInk}

\setbeamertemplate{navigation symbols}{}
\setbeamertemplate{footline}[frame number]
\setbeamertemplate{itemize item}{\color{CalPolyGold}$\blacktriangleright$}
\setbeamertemplate{itemize subitem}{\color{CalPolyGreen}$\bullet$}

\newcommand{\CourseNumber}{Course}
\newcommand{\CourseTitle}{Course Title}
\newcommand{\LectureNumber}{0}
\newcommand{\LectureTitle}{Lecture Title}
\newcommand{\LectureDate}{Date}
\newcommand{\CourseUnit}{Unit}

\newcommand{\coursenumber}[1]{\renewcommand{\CourseNumber}{#1}}
\newcommand{\coursetitle}[1]{\renewcommand{\CourseTitle}{#1}}
\newcommand{\lecturenumber}[1]{\renewcommand{\LectureNumber}{#1}}
\newcommand{\lecturetitle}[1]{\renewcommand{\LectureTitle}{#1}}
\newcommand{\lecturedate}[1]{\renewcommand{\LectureDate}{#1}}
\newcommand{\courseunit}[1]{\renewcommand{\CourseUnit}{#1}}

\author{Kirk Duran}
\institute{California Polytechnic State University, San Luis Obispo}
\date{\LectureDate}

\newcommand{\makedecktitle}{
  \begin{frame}[plain]
    \vspace{0.25in}
    {\usebeamerfont{title}\color{CalPolyGreen}\LectureTitle\par}
    \vspace{0.18in}
    {\large \CourseNumber{}: \CourseTitle\par}
    \vspace{0.08in}
    {\large Lecture \LectureNumber{} -- \LectureDate\par}
    \vfill
    {\small Kirk Duran \textbar{} Cal Poly, San Luis Obispo\par}
  \end{frame}
}

\newcommand{\placeholdernote}{
  \begin{block}{Placeholder}
    Replace these bullets with the final lecture material, examples, and in-class activities.
  \end{block}
}

\AtBeginSection[]{
  \begin{frame}{Roadmap}
    \tableofcontents[currentsection]
  \end{frame}
}


\pdfstringdefDisableCommands{\def\translate#1{#1}}

\graphicspath{{../assets/lecture-11/}}

% If an image file is not yet available, the slide displays a labeled
% placeholder using the intended filename. Place the matching file in
% ../assets/lecture-11/ to replace the placeholder automatically.
\newcommand{\lectureimage}[2][1.75in]{%
  \IfFileExists{../assets/lecture-11/#2}{%
    \includegraphics[width=\linewidth,height=#1,keepaspectratio]{#2}%
  }{%
    \fbox{%
      \parbox[c][#1][c]{0.94\linewidth}{%
        \centering
        \small\textbf{Image Placeholder}\\[0.08in]
        \scriptsize\texttt{\detokenize{#2}}
      }%
    }%
  }%
}

\newcommand{\lectureplaceholder}[2][1.75in]{%
  \fbox{%
    \parbox[c][#1][c]{0.94\linewidth}{%
      \centering
      \small\textbf{Image Placeholder}\\[0.08in]
      \scriptsize #2
    }%
  }%
}

\setbeamersize{text margin left=0.42in,text margin right=0.42in}

\coursenumber{CS 1001}
\coursetitle{Introduction to Computing}
\lecturenumber{11}
\lecturetitle{Scope and Call Frames}
\lecturedate{September 18, 2026}
\courseunit{1. Computing and Program Execution}

\begin{document}

\makedecktitle

\begin{frame}{Today}
  \begin{itemize}
    \item Define \textbf{local}, \textbf{enclosing}, and \textbf{global} scopes precisely.
    \item Separate lexical name resolution from runtime call-frame organization.
    \item Trace nested functions and identify free-variable bindings.
    \item Use \texttt{global} and \texttt{nonlocal} to change where assignment binds a name.
    \item Bind arguments using position, keywords, and default parameter values.
    \item Diagnose common scope and argument-binding errors by tracing program state.
  \end{itemize}

  \begin{keyidea}
    A name has meaning only relative to a scope, while a running function call has state only relative to its call frame.
  \end{keyidea}
\end{frame}

\sectionframe{Part 1: Scope and Name Resolution}

\begin{frame}{Names Are Resolved Within Scopes}
  \begin{columns}[T,onlytextwidth]
    \begin{column}{0.54\textwidth}
      \begin{cpdefinition}
        A \textbf{scope} is a region of program text in which a binding for a name is visible.
      \end{cpdefinition}

      \begin{itemize}
        \item Assignment, parameters, imports, and function definitions can bind names.
        \item Using a name requires Python to determine which visible binding the occurrence refers to.
        \item The nearest applicable lexical binding is selected.
      \end{itemize}
    \end{column}

    \begin{column}{0.40\textwidth}
      % Intended visual: nested lexical regions labeled local, enclosing, global, builtins.
      \lectureimage[2.55in]{scope-resolution-layers.png}
    \end{column}
  \end{columns}
\end{frame}

\begin{frame}{Local and Global Scope}
  \begin{columns}[T,onlytextwidth]
    \begin{column}{0.46\textwidth}
      \begin{cpdefinition}
        A \textbf{local binding} is a name bound within a function body and belonging to that function's local scope.
      \end{cpdefinition}

      \begin{block}{Example}
        \texttt{def f(x):}\\
        \hspace*{0.20in}\texttt{y = x + 1}
      \end{block}

      Here \texttt{x} and \texttt{y} are local to \texttt{f}.
    \end{column}

    \begin{column}{0.46\textwidth}
      \begin{cpdefinition}
        A \textbf{global binding} is a name bound at module level.
      \end{cpdefinition}

      \begin{block}{Example}
        \texttt{tax\_rate = 0.075}\\[0.08in]
        \texttt{def total(price):}\
        \hspace*{0.20in}\texttt{return price * (1 + tax\_rate)}
      \end{block}
    \end{column}
  \end{columns}
\end{frame}

\begin{frame}{A Function May Read a Global Binding Without \texttt{global}}
  \begin{block}{Program}
    \texttt{rate = 2}\\[0.06in]
    \texttt{def scale(value):}\\
    \hspace*{0.20in}\texttt{return value * rate}\\[0.08in]
    \texttt{answer = scale(5)}
  \end{block}

  \begin{enumerate}
    \item \texttt{value} is found in the local scope of \texttt{scale}.
    \item \texttt{rate} has no local binding in \texttt{scale}.
    \item Name resolution continues outward and finds global \texttt{rate -> 2}.
    \item The function returns \texttt{10}.
  \end{enumerate}

  \begin{keyidea}
    Reading a global name and rebinding a global name are different operations.
  \end{keyidea}
\end{frame}

\begin{frame}{Assignment Inside a Function Normally Creates a Local Binding}
  \begin{block}{Program}
    \texttt{count = 10}\\[0.06in]
    \texttt{def reset():}\\
    \hspace*{0.20in}\texttt{count = 0}\\
    \hspace*{0.20in}\texttt{return count}\\[0.08in]
    \texttt{x = reset()}
  \end{block}

  \begin{columns}[T,onlytextwidth]
    \begin{column}{0.46\textwidth}
      \begin{block}{Global binding}
        \texttt{count -> 10}
      \end{block}
    \end{column}

    \begin{column}{0.46\textwidth}
      \begin{block}{During \texttt{reset}}
        local \texttt{count -> 0}\\
        return \texttt{0}
      \end{block}
    \end{column}
  \end{columns}

  The module-level \texttt{count} remains \texttt{10}.
\end{frame}

\begin{frame}{Shadowing Creates a Different Binding with the Same Name}
  \begin{cpdefinition}
    \textbf{Shadowing} occurs when a nearer scope binds a name that also exists in an outer scope.
  \end{cpdefinition}

  \begin{block}{Example}
    \texttt{price = 100}\\[0.06in]
    \texttt{def discounted():}\\
    \hspace*{0.20in}\texttt{price = 80}\\
    \hspace*{0.20in}\texttt{return price}
  \end{block}

  \begin{itemize}
    \item Inside \texttt{discounted}, \texttt{price} resolves to the local binding \texttt{80}.
    \item Outside the function, \texttt{price} still resolves to the global binding \texttt{100}.
    \item Same spelling does not imply same binding.
  \end{itemize}
\end{frame}

\begin{frame}{A Subtle Rule: Binding Anywhere Makes the Name Local}
  \begin{block}{Program}
    \texttt{x = 10}\\[0.06in]
    \texttt{def f():}\\
    \hspace*{0.20in}\texttt{print(x)}\\
    \hspace*{0.20in}\texttt{x = 20}
  \end{block}

  \begin{alertblock}{Result}
    Calling \texttt{f()} raises \texttt{UnboundLocalError}.
  \end{alertblock}

  \begin{itemize}
    \item Because \texttt{x} is assigned somewhere in \texttt{f}, Python treats \texttt{x} as local throughout that function block.
    \item The first \texttt{print(x)} therefore attempts to read the local \texttt{x} before that local binding has received a value.
  \end{itemize}

  \begin{keyidea}
    Scope classification is not determined statement-by-statement during execution.
  \end{keyidea}
\end{frame}

\sectionframe{Part 2: Call Frames as Runtime State}

\begin{frame}{A Call Frame Represents One Active Execution of a Function}
  \begin{cpdefinition}
    A \textbf{call frame} is the runtime execution state associated with one active function call.
  \end{cpdefinition}

  \begin{columns}[T,onlytextwidth]
    \begin{column}{0.52\textwidth}
      A useful CS 1001 model records:
      \begin{itemize}
        \item parameter bindings;
        \item local-variable bindings;
        \item the currently executing location;
        \item the relationship to the caller.
      \end{itemize}
    \end{column}

    \begin{column}{0.42\textwidth}
      % Intended visual: one frame with parameters, locals, instruction location, caller link.
      \lectureimage[2.55in]{call-frame-anatomy.png}
    \end{column}
  \end{columns}
\end{frame}

\begin{frame}{Frame Lifetime Follows Call Lifetime}
  \begin{enumerate}
    \item Evaluate the arguments at the call site.
    \item Begin the call and initialize the frame's parameter bindings.
    \item Execute the function body, adding or changing local bindings.
    \item Evaluate a \texttt{return} expression or reach the end of the function.
    \item End the active call; the caller resumes using the produced value.
  \end{enumerate}

  \begin{columns}[T,onlytextwidth]
    \begin{column}{0.45\textwidth}
      \begin{block}{During call}
        frame exists\\
        locals are inspectable
      \end{block}
    \end{column}

    \begin{column}{0.45\textwidth}
      \begin{block}{After return}
        call is no longer active\\
        its ordinary local environment is no longer active
      \end{block}
    \end{column}
  \end{columns}
\end{frame}

\begin{frame}{Two Calls to the Same Function Use Different Frames}
  \begin{block}{Program}
    \texttt{def increase(value):}\\
    \hspace*{0.20in}\texttt{result = value + 1}\\
    \hspace*{0.20in}\texttt{return result}\\[0.08in]
    \texttt{a = increase(4)}\\
    \texttt{b = increase(10)}
  \end{block}

  \begin{columns}[T,onlytextwidth]
    \begin{column}{0.45\textwidth}
      \begin{block}{First call}
        \texttt{value -> 4}\\
        \texttt{result -> 5}
      \end{block}
    \end{column}

    \begin{column}{0.45\textwidth}
      \begin{block}{Second call}
        \texttt{value -> 10}\\
        \texttt{result -> 11}
      \end{block}
    \end{column}
  \end{columns}

  The parameter name \texttt{value} is reused, but the two runtime bindings belong to different calls.
\end{frame}

\begin{frame}{Call Stack and Scope Answer Different Questions}
  \begin{columns}[T,onlytextwidth]
    \begin{column}{0.46\textwidth}
      \begin{block}{Call stack}
        Runtime structure.\\[0.08in]
        \textbf{Question:} Which calls are active, and which caller resumes next?
      \end{block}
    \end{column}

    \begin{column}{0.46\textwidth}
      \begin{block}{Lexical scope}
        Program-text structure.\\[0.08in]
        \textbf{Question:} Which binding does this occurrence of a name refer to?
      \end{block}
    \end{column}
  \end{columns}

  \begin{alertblock}{Important}
    A caller's local variables do not automatically become visible to a callee merely because the caller's frame is beneath it on the call stack.
  \end{alertblock}
\end{frame}

\begin{frame}{Trace a Call Frame, Not Just the Final Result}
  \begin{block}{Program}
    \texttt{def compute(a, b):}\\
    \hspace*{0.20in}\texttt{total = a + b}\\
    \hspace*{0.20in}\texttt{scaled = total * 2}\\
    \hspace*{0.20in}\texttt{return scaled}\\[0.08in]
    \texttt{answer = compute(3, 5)}
  \end{block}

  \begin{center}
    \renewcommand{\arraystretch}{1.25}
    \begin{tabular}{l|l}
      \textbf{Execution point} & \textbf{Local frame bindings} \\\hline
      function entry & \texttt{a -> 3, b -> 5} \\
      after \texttt{total = a + b} & \texttt{a -> 3, b -> 5, total -> 8} \\
      after \texttt{scaled = total * 2} & add \texttt{scaled -> 16} \\
      return & produce \texttt{16}
    \end{tabular}
  \end{center}
\end{frame}

\sectionframe{Part 3: Nested Functions and Enclosing Scope}

\begin{frame}{A Function Definition May Appear Inside Another Function}
  \begin{block}{Example}
    \texttt{def outer(value):}\\
    \hspace*{0.20in}\texttt{def double(x):}\\
    \hspace*{0.40in}\texttt{return x * 2}\\
    \hspace*{0.20in}\texttt{return double(value)}
  \end{block}

  \begin{itemize}
    \item Executing the inner \texttt{def} statement creates a function object and binds it to the local name \texttt{double}.
    \item That binding belongs to the current call of \texttt{outer}.
    \item The inner function can be called like any other function value while its binding is available.
  \end{itemize}
\end{frame}

\begin{frame}{Nested Functions Add an Enclosing Function Scope}
  \begin{block}{Program}
    \texttt{def outer(base):}\\
    \hspace*{0.20in}\texttt{offset = 3}\\
    \hspace*{0.20in}\texttt{def inner(x):}\\
    \hspace*{0.40in}\texttt{return x + offset}\\
    \hspace*{0.20in}\texttt{return inner(base)}
  \end{block}

  \begin{itemize}
    \item \texttt{x} is local to \texttt{inner}.
    \item \texttt{offset} is not local to \texttt{inner}; it is a \textbf{free variable} there.
    \item Python resolves \texttt{offset} in the nearest enclosing function scope where it is bound.
  \end{itemize}

  \begin{keyidea}
    Lexical nesting, not the order of callers on the stack, determines enclosing-scope name resolution.
  \end{keyidea}
\end{frame}

\begin{frame}{Trace a Nested Function Call}
  \begin{block}{Call}
    \centering
    \Large\texttt{result = outer(5)}
  \end{block}

  \begin{columns}[T,onlytextwidth]
    \begin{column}{0.46\textwidth}
      \begin{block}{\texttt{outer} frame}
        \texttt{base -> 5}\\
        \texttt{offset -> 3}\\
        \texttt{inner -> <function>}
      \end{block}
    \end{column}

    \begin{column}{0.46\textwidth}
      \begin{block}{\texttt{inner} frame}
        \texttt{x -> 5}\\
        free \texttt{offset -> 3} from enclosing scope\\
        return \texttt{8}
      \end{block}
    \end{column}
  \end{columns}

  Then \texttt{outer} returns \texttt{8}, so \texttt{result -> 8}.
\end{frame}

\begin{frame}{\texttt{nonlocal} Rebinds the Nearest Enclosing Function Binding}
  \begin{block}{Program}
    \texttt{def outer():}\\
    \hspace*{0.20in}\texttt{count = 0}\\
    \hspace*{0.20in}\texttt{def step():}\\
    \hspace*{0.40in}\texttt{nonlocal count}\\
    \hspace*{0.40in}\texttt{count = count + 1}\\
    \hspace*{0.20in}\texttt{step()}\\
    \hspace*{0.20in}\texttt{return count}
  \end{block}

  \begin{itemize}
    \item \texttt{nonlocal count} prevents assignment from creating a new local \texttt{count} in \texttt{step}.
    \item The assignment targets the existing \texttt{count} in the nearest enclosing function scope.
    \item \texttt{outer()} therefore returns \texttt{1}.
  \end{itemize}
\end{frame}

\sectionframe{Part 4: The \texttt{global} Statement}

\begin{frame}{\texttt{global} Changes the Binding Target of a Name}
  \begin{block}{Program}
    \texttt{count = 0}\\[0.06in]
    \texttt{def increment():}\\
    \hspace*{0.20in}\texttt{global count}\\
    \hspace*{0.20in}\texttt{count = count + 1}
  \end{block}

  \begin{cpdefinition}
    Within the current scope, \texttt{global name} declares that uses and assignments of \texttt{name} refer to the module-level global binding rather than to a new local binding.
  \end{cpdefinition}

  \begin{keyidea}
    \texttt{global} is a lexical declaration about name resolution and assignment target selection; it does not copy a global value into the function. It must precede uses of that name in the same function block.
  \end{keyidea}
\end{frame}

\begin{frame}{Trace a Global Rebinding}
  \begin{block}{Program}
    \texttt{score = 10}\\[0.06in]
    \texttt{def add\_points(points):}\\
    \hspace*{0.20in}\texttt{global score}\\
    \hspace*{0.20in}\texttt{score = score + points}\\
    \hspace*{0.20in}\texttt{return score}\\[0.08in]
    \texttt{result = add\_points(5)}
  \end{block}

  \begin{enumerate}
    \item Frame begins with local parameter \texttt{points -> 5}.
    \item \texttt{score} resolves to global \texttt{10} because of the declaration.
    \item Assignment rebinds global \texttt{score -> 15}.
    \item The call returns \texttt{15}; then \texttt{result -> 15} globally.
  \end{enumerate}
\end{frame}

\begin{frame}{Without \texttt{global}, the Same Code Is Not Equivalent}
  \begin{block}{Problematic version}
    \texttt{score = 10}\\[0.06in]
    \texttt{def add\_points(points):}\\
    \hspace*{0.20in}\texttt{score = score + points}\\
    \hspace*{0.20in}\texttt{return score}
  \end{block}

  \begin{alertblock}{Why it fails}
    The assignment makes \texttt{score} local to \texttt{add\_points}. The right-hand side then tries to read that local \texttt{score} before it has been bound, producing \texttt{UnboundLocalError}.
  \end{alertblock}

  \begin{keyidea}
    Adding an assignment can change how every occurrence of that name is classified within the function block.
  \end{keyidea}
\end{frame}

\begin{frame}{Local, \texttt{nonlocal}, and \texttt{global} Select Different Binding Targets}
  \begin{center}
    \renewcommand{\arraystretch}{1.35}
    \begin{tabular}{l|l|l}
      \textbf{Inside a function} & \textbf{Assignment target} & \textbf{Typical meaning} \\\hline
      ordinary assignment & current local scope & create/rebind local name \\
      \texttt{nonlocal x} & nearest enclosing function binding & rebind enclosing state \\
      \texttt{global x} & module-level global binding & rebind global state
    \end{tabular}
  \end{center}

  \begin{keyidea}
    These declarations alter where a name is resolved for assignment; they do not change the right-hand-side evaluation rules. Prefer explicit parameters and return values when shared mutable state is unnecessary.
  \end{keyidea}
\end{frame}

\sectionframe{Part 5: Parameter Binding Beyond Position}

\begin{frame}{Positional Arguments Bind by Parameter Position}
  \begin{block}{Definition and call}
    \texttt{def rectangle(width, height):}\
    \hspace*{0.20in}\texttt{return width * height}\\[0.08in]
    \texttt{area = rectangle(5, 3)}
  \end{block}

  \begin{block}{Binding}
    first argument \texttt{5 -> width}\\
    second argument \texttt{3 -> height}
  \end{block}

  Position determines the correspondence; the caller's variable names, if any, are irrelevant after their values are evaluated.
\end{frame}

\begin{frame}{Keyword Arguments Bind by Parameter Name}
  \begin{block}{Calls}
    \texttt{rectangle(width=5, height=3)}\\
    \texttt{rectangle(height=3, width=5)}
  \end{block}

  \begin{itemize}
    \item A keyword argument explicitly identifies the parameter slot to fill.
    \item Keyword argument order therefore does not determine parameter correspondence.
    \item The argument expression is still evaluated to obtain a value before the call executes.
  \end{itemize}

  \begin{keyidea}
    Positional binding answers ``which slot by order?''; keyword binding answers ``which slot by parameter name?''
  \end{keyidea}
\end{frame}

\begin{frame}{Positional and Keyword Arguments May Be Mixed}
  \begin{block}{Definition}
    \texttt{def describe(name, age, city):}\
    \hspace*{0.20in}\texttt{return name, age, city}
  \end{block}

  \begin{block}{Valid call}
    \texttt{describe("Ada", age=20, city="SLO")}
  \end{block}

  \begin{enumerate}
    \item Positional argument \texttt{"Ada"} fills \texttt{name}.
    \item Keyword \texttt{age=20} fills \texttt{age}.
    \item Keyword \texttt{city="SLO"} fills \texttt{city}.
  \end{enumerate}

  In ordinary calls, explicit positional arguments precede explicit keyword arguments.

  \begin{alertblock}{Binding must be unambiguous}
    A parameter may not receive two values, and every required parameter must be filled. Duplicate, missing, or unknown arguments raise \texttt{TypeError}.
  \end{alertblock}
\end{frame}

\begin{frame}{Default Parameter Values Fill Omitted Slots}
  \begin{block}{Definition}
    \texttt{def power(base, exponent=2):}\
    \hspace*{0.20in}\texttt{return base ** exponent}
  \end{block}

  \begin{columns}[T,onlytextwidth]
    \begin{column}{0.45\textwidth}
      \begin{block}{\texttt{power(5)}}
        \texttt{base -> 5}\\
        \texttt{exponent -> 2} from default\\
        return \texttt{25}
      \end{block}
    \end{column}

    \begin{column}{0.45\textwidth}
      \begin{block}{\texttt{power(5, exponent=3)}}
        \texttt{base -> 5}\\
        \texttt{exponent -> 3} from argument\\
        return \texttt{125}
      \end{block}
    \end{column}
  \end{columns}
\end{frame}

\begin{frame}{Default Expressions Are Evaluated When the Function Is Defined}
  \begin{block}{Trace carefully}
    \texttt{x = 10}\\
    \texttt{def show(value=x):}\\
    \hspace*{0.20in}\texttt{return value}\\[0.06in]
    \texttt{x = 20}\\
    \texttt{answer = show()}
  \end{block}

  \begin{enumerate}
    \item When the \texttt{def} executes, the default expression \texttt{x} evaluates to \texttt{10}.
    \item Rebinding global \texttt{x -> 20} does not recompute the stored default.
    \item \texttt{show()} therefore binds \texttt{value -> 10}.
  \end{enumerate}

  \begin{alertblock}{Design consequence}
    Mutable objects used as defaults are shared across calls that omit that argument; this is usually avoided in introductory code.
  \end{alertblock}
\end{frame}

\begin{frame}{Python Supports Multiple Parameter Kinds}
  \begin{block}{General form}
    \centering
    \texttt{def f(a, b, /, c, d=0, *, scale=1): ...}
  \end{block}

  \begin{center}
    \renewcommand{\arraystretch}{1.25}
    \begin{tabular}{l|l}
      \textbf{Parameters} & \textbf{How arguments may be supplied} \\\hline
      \texttt{a, b} before \texttt{/} & positional only \\
      \texttt{c, d} & positional or keyword \\
      \texttt{scale} after \texttt{*} & keyword only
    \end{tabular}
  \end{center}

  \begin{keyidea}
    Most student-defined functions use positional-or-keyword parameters; \texttt{/} and \texttt{*} make an interface more explicit when needed.
  \end{keyidea}
\end{frame}

\sectionframe{Part 6: Integrated Tracing and Debugging}

\begin{frame}{Integrated Trace: Global and Local Bindings with Keywords}
  \begin{block}{Program}
    \texttt{rate = 2}\\
    \texttt{def transform(value, offset=1):}\\
    \hspace*{0.20in}\texttt{scaled = value * rate}\\
    \hspace*{0.20in}\texttt{return scaled + offset}\\[0.08in]
    \texttt{answer = transform(offset=4, value=3)}
  \end{block}

  \begin{enumerate}
    \item Keyword binding creates local \texttt{value -> 3}, \texttt{offset -> 4}.
    \item Global lookup supplies \texttt{rate -> 2}.
    \item Local assignment creates \texttt{scaled -> 6}.
    \item Return expression evaluates \texttt{6 + 4 -> 10}.
    \item The caller binds global \texttt{answer -> 10}.
  \end{enumerate}
\end{frame}

\begin{frame}{Integrated Trace: Nested Function and Enclosing Binding}
  \begin{block}{Program}
    \texttt{def make\_total(base, tax=0.10):}\\
    \hspace*{0.20in}\texttt{def add\_tax(amount):}\\
    \hspace*{0.40in}\texttt{return amount * (1 + tax)}\\
    \hspace*{0.20in}\texttt{return add\_tax(base)}\\[0.08in]
    \texttt{result = make\_total(50, tax=0.20)}
  \end{block}

  \begin{columns}[T,onlytextwidth]
    \begin{column}{0.45\textwidth}
      \begin{block}{Outer frame}
        \texttt{base -> 50}\\
        \texttt{tax -> 0.20}\\
        \texttt{add\_tax -> <function>}
      \end{block}
    \end{column}
    \begin{column}{0.45\textwidth}
      \begin{block}{Inner frame}
        \texttt{amount -> 50}\\
        free \texttt{tax -> 0.20}\\
        return \texttt{60.0}
      \end{block}
    \end{column}
  \end{columns}
\end{frame}

\begin{frame}{Debugger Strategy: Inspect the Correct Frame}
  \begin{columns}[T,onlytextwidth]
    \begin{column}{0.54\textwidth}
      When stopped inside nested calls:
      \begin{enumerate}
        \item Identify the currently executing frame.
        \item Inspect its parameter and local bindings.
        \item Use the call stack to select an outer active frame when needed.
        \item Keep lexical scope separate from the debugger's stack ordering.
        \item Predict name resolution before relying on the debugger display.
      \end{enumerate}
    \end{column}

    \begin{column}{0.40\textwidth}
      % Intended visual: debugger call-stack panel with two frames and separate locals panels.
      \lectureimage[2.70in]{debugger-frame-selection.png}
    \end{column}
  \end{columns}
\end{frame}

\begin{frame}{Tracing Activity: Resolve Every Name Before Computing}
  \begin{block}{Predict without running}
    \texttt{x = 5}\\[0.04in]
    \texttt{def outer(x, scale=2):}\\
    \hspace*{0.20in}\texttt{offset = 3}\\
    \hspace*{0.20in}\texttt{def inner(value):}\\
    \hspace*{0.40in}\texttt{return value * scale + offset}\\
    \hspace*{0.20in}\texttt{return inner(x)}\\[0.06in]
    \texttt{answer = outer(scale=4, x=2)}
  \end{block}

  \begin{activity}
    Draw the active frames during \texttt{inner}. For each occurrence of \texttt{x}, \texttt{scale}, \texttt{offset}, and \texttt{value}, identify its binding and scope. Then determine \texttt{answer}.
  \end{activity}
\end{frame}

\begin{frame}{Takeaways}
  \begin{itemize}
    \item \textbf{Scope} determines which binding a name occurrence can refer to; a \textbf{call frame} records runtime state for one active call.
    \item Ordinary assignment inside a function creates or rebinds a local name unless \texttt{nonlocal} or \texttt{global} changes the binding target.
    \item Nested functions may read names from lexically enclosing function scopes.
    \item \texttt{global} permits rebinding module-level names, but it also introduces externally visible state changes that must be traced explicitly.
    \item Function-call argument binding fills parameter slots using positional arguments, keyword arguments, and default values.
    \item Accurate tracing separates three questions: \emph{Which frame is active? Which scope supplies this name? Which value is bound?}
  \end{itemize}

  \begin{keyidea}
    The central discipline is to locate the binding first, then evaluate the expression.
  \end{keyidea}
\end{frame}

\end{document}
